% Options for packages loaded elsewhere
\PassOptionsToPackage{unicode}{hyperref}
\PassOptionsToPackage{hyphens}{url}
%
\documentclass[
  11pt,
]{article}
\usepackage{amsmath,amssymb}
\usepackage{setspace}
\usepackage{iftex}
\ifPDFTeX
  \usepackage[T1]{fontenc}
  \usepackage[utf8]{inputenc}
  \usepackage{textcomp} % provide euro and other symbols
\else % if luatex or xetex
  \usepackage{unicode-math} % this also loads fontspec
  \defaultfontfeatures{Scale=MatchLowercase}
  \defaultfontfeatures[\rmfamily]{Ligatures=TeX,Scale=1}
\fi
\usepackage{lmodern}
\ifPDFTeX\else
  % xetex/luatex font selection
\fi
% Use upquote if available, for straight quotes in verbatim environments
\IfFileExists{upquote.sty}{\usepackage{upquote}}{}
\IfFileExists{microtype.sty}{% use microtype if available
  \usepackage[]{microtype}
  \UseMicrotypeSet[protrusion]{basicmath} % disable protrusion for tt fonts
}{}
\makeatletter
\@ifundefined{KOMAClassName}{% if non-KOMA class
  \IfFileExists{parskip.sty}{%
    \usepackage{parskip}
  }{% else
    \setlength{\parindent}{0pt}
    \setlength{\parskip}{6pt plus 2pt minus 1pt}}
}{% if KOMA class
  \KOMAoptions{parskip=half}}
\makeatother
\usepackage{xcolor}
\usepackage[margin=1in]{geometry}
\usepackage{color}
\usepackage{fancyvrb}
\newcommand{\VerbBar}{|}
\newcommand{\VERB}{\Verb[commandchars=\\\{\}]}
\DefineVerbatimEnvironment{Highlighting}{Verbatim}{commandchars=\\\{\}}
% Add ',fontsize=\small' for more characters per line
\newenvironment{Shaded}{}{}
\newcommand{\AlertTok}[1]{\textcolor[rgb]{1.00,0.00,0.00}{\textbf{#1}}}
\newcommand{\AnnotationTok}[1]{\textcolor[rgb]{0.38,0.63,0.69}{\textbf{\textit{#1}}}}
\newcommand{\AttributeTok}[1]{\textcolor[rgb]{0.49,0.56,0.16}{#1}}
\newcommand{\BaseNTok}[1]{\textcolor[rgb]{0.25,0.63,0.44}{#1}}
\newcommand{\BuiltInTok}[1]{\textcolor[rgb]{0.00,0.50,0.00}{#1}}
\newcommand{\CharTok}[1]{\textcolor[rgb]{0.25,0.44,0.63}{#1}}
\newcommand{\CommentTok}[1]{\textcolor[rgb]{0.38,0.63,0.69}{\textit{#1}}}
\newcommand{\CommentVarTok}[1]{\textcolor[rgb]{0.38,0.63,0.69}{\textbf{\textit{#1}}}}
\newcommand{\ConstantTok}[1]{\textcolor[rgb]{0.53,0.00,0.00}{#1}}
\newcommand{\ControlFlowTok}[1]{\textcolor[rgb]{0.00,0.44,0.13}{\textbf{#1}}}
\newcommand{\DataTypeTok}[1]{\textcolor[rgb]{0.56,0.13,0.00}{#1}}
\newcommand{\DecValTok}[1]{\textcolor[rgb]{0.25,0.63,0.44}{#1}}
\newcommand{\DocumentationTok}[1]{\textcolor[rgb]{0.73,0.13,0.13}{\textit{#1}}}
\newcommand{\ErrorTok}[1]{\textcolor[rgb]{1.00,0.00,0.00}{\textbf{#1}}}
\newcommand{\ExtensionTok}[1]{#1}
\newcommand{\FloatTok}[1]{\textcolor[rgb]{0.25,0.63,0.44}{#1}}
\newcommand{\FunctionTok}[1]{\textcolor[rgb]{0.02,0.16,0.49}{#1}}
\newcommand{\ImportTok}[1]{\textcolor[rgb]{0.00,0.50,0.00}{\textbf{#1}}}
\newcommand{\InformationTok}[1]{\textcolor[rgb]{0.38,0.63,0.69}{\textbf{\textit{#1}}}}
\newcommand{\KeywordTok}[1]{\textcolor[rgb]{0.00,0.44,0.13}{\textbf{#1}}}
\newcommand{\NormalTok}[1]{#1}
\newcommand{\OperatorTok}[1]{\textcolor[rgb]{0.40,0.40,0.40}{#1}}
\newcommand{\OtherTok}[1]{\textcolor[rgb]{0.00,0.44,0.13}{#1}}
\newcommand{\PreprocessorTok}[1]{\textcolor[rgb]{0.74,0.48,0.00}{#1}}
\newcommand{\RegionMarkerTok}[1]{#1}
\newcommand{\SpecialCharTok}[1]{\textcolor[rgb]{0.25,0.44,0.63}{#1}}
\newcommand{\SpecialStringTok}[1]{\textcolor[rgb]{0.73,0.40,0.53}{#1}}
\newcommand{\StringTok}[1]{\textcolor[rgb]{0.25,0.44,0.63}{#1}}
\newcommand{\VariableTok}[1]{\textcolor[rgb]{0.10,0.09,0.49}{#1}}
\newcommand{\VerbatimStringTok}[1]{\textcolor[rgb]{0.25,0.44,0.63}{#1}}
\newcommand{\WarningTok}[1]{\textcolor[rgb]{0.38,0.63,0.69}{\textbf{\textit{#1}}}}
\setlength{\emergencystretch}{3em} % prevent overfull lines
\providecommand{\tightlist}{%
  \setlength{\itemsep}{0pt}\setlength{\parskip}{0pt}}
\setcounter{secnumdepth}{-\maxdimen} % remove section numbering
\usepackage{tikz}
\ifLuaTeX
  \usepackage{selnolig}  % disable illegal ligatures
\fi
\usepackage{bookmark}
\IfFileExists{xurl.sty}{\usepackage{xurl}}{} % add URL line breaks if available
\urlstyle{same}
\hypersetup{
  pdftitle={Research That Can Continue},
  pdfauthor={Watson Hartsoe},
  pdfkeywords={research objects, provenance, recursive inquiry, language
models, digital preservation, knowledge graphs, reproducible research},
  hidelinks,
  pdfcreator={LaTeX via pandoc}}

\title{Research That Can Continue}
\usepackage{etoolbox}
\makeatletter
\providecommand{\subtitle}[1]{% add subtitle to \maketitle
  \apptocmd{\@title}{\par {\large #1 \par}}{}{}
}
\makeatother
\subtitle{Recursive Objects, Typed Absence, and the Checkpoint as a
Scholarly Form}
\author{Watson Hartsoe}
\date{18 August 2026}

\begin{document}
\maketitle
\begin{abstract}
Research conducted with language models is unusually vulnerable to a
mundane form of loss: the conversation disappears while the research
remains unfinished. Saving the prose is not enough. What disappears with
the interface is the structure of continuation - which question remained
live, which source was still missing, which contradiction mattered,
which test could discriminate rival readings, and which interpretive
layer could be discarded without destroying evidence. This working paper
develops the idea of \emph{continuable scholarship}: a research object
should preserve not only what has been said but enough epistemic and
operational state for another researcher or model to make the next
warranted move. I examine a four-layer architecture consisting of atomic
source-grounded zettels, type-invariant recursive forage, graph-level
allocation of research attention, and portable checkpoints. The central
argument is that durability requires different temporal rules for
evidence and interpretation: evidence should accumulate monotonically,
while interpretation must remain nonmonotonic and replaceable. From this
follow three further requirements: unresolved addresses may steer
inquiry without becoming evidence; production provenance must be
separated from evidentiary support; and deterministic verification
should be delegated to tools rather than rhetorical confidence. The
result is less a note-taking system than a proposal for a scholarly
object that can come back alive after its originating context is gone.
\end{abstract}

\setstretch{1.12}
\begin{tikzpicture}[remember picture,overlay]
  \coordinate (m) at ([xshift=-0.62in,yshift=0.58in]current page.south east);
  \fill (m) circle (0.8pt);
  \draw[line width=0.35pt] ([xshift=-2pt,yshift=5pt]m) arc[start angle=220,end angle=320,radius=7pt];
  \draw[line width=0.35pt] ([xshift=-9pt,yshift=7pt]m) arc[start angle=150,end angle=255,radius=12pt];
  \node[anchor=south east,align=right,font=\fontsize{6.5}{8}\selectfont,text=black!65] at ([xshift=-18pt,yshift=-1pt]m) {Working Paper $\cdot$ AI-augmented research process\\Evidence, prompts, and making history preserved};
\end{tikzpicture}

\section{1. The missing object is not the
note}\label{the-missing-object-is-not-the-note}

The ordinary description of research preservation starts too late. It
asks whether a document can still be opened, whether its bytes remain
intact, whether its citations resolve, or whether its workflow can be
reproduced. Those are necessary questions, but research often becomes
unusable before any of them fail. A note can remain perfectly readable
while the inquiry that produced it has disappeared.

This is especially visible in research conducted through language-model
conversations. A working session may contain a source passage, a
provisional interpretation, a correction, an unresolved disagreement, a
promising citation, and a proposed test. If the session is reduced to a
polished summary, the conclusions may survive while the \emph{handles
for continuing the work} vanish. The loss is not semantic content alone.
It is the loss of research state.

The SLIPCASE / FORAGE architecture begins from a different preservation
question: what has to be true of a research object for another
researcher or capable model to continue the inquiry without the
originating conversation? The first zettel in the present field names
the criterion as \textbf{continuability}: a durable object must carry
enough identity, evidence, relational position, unresolvedness, and
possible action to permit a warranted next move (Z-SLIP-001). That
proposal shifts the unit of preservation. The object to be preserved is
not a statement but a \emph{statement-in-an-inquiry}.

This differs from established packaging and provenance work by emphasis
rather than opposition. BagIt defines a minimal filesystem convention
for reliable storage and transfer of arbitrary digital content {[}1{]}.
W3C PROV provides a general model for representing entities, activities,
agents, and derivations so provenance can be exchanged across systems
{[}2{]}. RO-Crate packages research artifacts with machine-readable
metadata for reuse, preservation, and reproducibility {[}3{]}. These
systems establish that a research package should be inspectable,
attributable, and transferable. The present wager is narrower and more
operational: a package should also expose enough unfinished structure to
make \emph{the next research operation} recoverable.

That requires distinguishing preservation from continuation. A complete
archive of yesterday can still be a poor instrument for tomorrow.

\section{2. The recursive object}\label{the-recursive-object}

The first layer is deliberately small: one atomic zettel. Its unusual
property is not atomism by itself but \textbf{type invariance}.
Recursive forage is formally described as:

\begin{Shaded}
\begin{Highlighting}[]
\NormalTok{FORAGE(ZETTEL) {-}\textgreater{} ZETTEL[]}
\end{Highlighting}
\end{Shaded}

Every child preserves the same field names and field order as its
parent. No child-only \texttt{PARENT}, \texttt{STATUS},
\texttt{ANSWERS}, or workflow metadata is permitted. A child must be
able to become a parent without rewriting.

This constraint turns a formatting choice into an architectural
boundary. Systems that append workflow-specific metadata at every
generation gradually make their outputs dependent on the tool that
produced them. The zettel instead carries only research-state fields:
source, passage, research object, question, mechanism, tension, missing
element, boundary, citation trail, test, platform, links, and
bibliography. Lineage uses the same representational vocabulary
available to every other card. The object remains recursively operable
because it does not acquire a new ontological status merely by being
produced later.

The resulting card has a peculiar temporal structure. It records what
has been learned while deliberately exposing what remains unresolved.
\texttt{QUESTION}, \texttt{TENSION}, \texttt{MISSING},
\texttt{BOUNDARY}, \texttt{CITATION\ TRAIL}, and \texttt{TEST} are not
metadata around the finding; they are part of the finding because they
delimit how the evidence can be used and where inquiry can proceed.

This explains why a card can be readable yet dead. If one preserves only
its thesis, the next operation must be invented from scratch. If one
preserves the strongest unresolved edge and a discriminating test, the
object can be resumed.

Yet recursion reveals a limit. Z-SLIP-013 notes that a self-forageable
card is not a self-directing research system. Once there are many cards,
a second state becomes necessary: which edges have already been
investigated, which contradictions remain open, which sources recur
across branches, and which possible next move is worth spending
attention on? That state should not be stuffed back into the zettel,
because doing so would destroy type invariance. The research object and
the research controller therefore separate.

\section{3. The frontier: genealogy is not
attention}\label{the-frontier-genealogy-is-not-attention}

Recursive systems tempt depth-first continuation: produce a child, then
forage the child, then forage its child. That preserves genealogy but
confuses descent with priority.

The autonomous graph layer makes a different claim. Parentage tells us
where a question came from. It does not tell us whether that question
deserves the next hour. The controller therefore surveys an
\textbf{active frontier} consisting both of local open fields and
relations between cards: contradictions, repeated vocabulary, competing
genealogies, source/implementation mismatches, unexpected crossings, or
shared sources that connect distant branches.

Its governing question is not ``what follows from the latest card?'' but
``if I could know one new thing next, which finding would most change
the graph?'' Z-SLIP-014 interprets this as a theory of
research-attention allocation.

The gain is obvious: a minor leaf no longer monopolizes research merely
because it was generated last. The danger is equally important. Any
estimate of ``expected epistemic gain'' is computed from features
already represented in the graph. The archive therefore risks rewarding
what it already knows how to name. A low-frequency anomaly, an
unfamiliar vocabulary, or a source outside the current conceptual
language may be more transformative precisely because the existing graph
cannot score it well.

The architecture partly protects itself by treating exploration as a
distinct mode and by periodically ``looking sideways'' for collisions.
But the deeper issue remains open: a research system that directs its
own attention needs a theory of \emph{representation blindness}.
Optimization over the frontier cannot discover what has left no trace in
the frontier.

This is one reason human deference is not merely a safety stop.
Z-SLIP-018 identifies a stronger condition: automatic inquiry should
yield when the remaining problem is not ``which evidence would
discriminate these readings?'' but ``which question ought to matter?''
The unresolved technical problem is how a system can detect that
boundary without already possessing the normative judgment it is
supposed to defer.

\section{4. Evidence should accumulate; interpretations should be
allowed to
die}\label{evidence-should-accumulate-interpretations-should-be-allowed-to-die}

The most consequential design move appears in the separation of three
layers:

\begin{Shaded}
\begin{Highlighting}[]
\NormalTok{EVIDENCE       ZETTEL · SOURCE · RESOURCE · PROMPT · APPEARANCE}
\NormalTok{FIELD          PLATFORM · LINK · CONCEPT · GHOST · BACKLINK · LINEAGE}
\NormalTok{INTERPRETATION MOC · ARRANGEMENT · TRAIL · PAPER}
\end{Highlighting}
\end{Shaded}

The layers obey different rules. Evidence is preserved. Field structure
is compiled. Interpretation remains revisable.

Z-SLIP-015 makes the temporal consequence explicit. The evidentiary
history should be approximately monotonic: new objects may be appended,
but previously preserved payloads are not silently rewritten to fit
later theory. Interpretation is nonmonotonic: a taxonomy may fail, a
genealogy may collapse, a graph bridge may prove superficial, and a
paper thesis may need to be abandoned.

This is not a claim that preserved evidence is interpretation-free.
Admission, segmentation, naming, and source selection are already
judgments. The claim is more practical: a later change of mind should
not require retroactively editing the historical object that made the
change of mind possible.

That principle changes the status of the paper. In ordinary scholarly
workflow, notes are often treated as disposable scaffolding for the
stable publication. Here the relation is reversed. The cards persist;
the paper is a compiled view over them. Z-SLIP-007 calls the paper ``a
wager'': one defensible traversal of a field that may later support a
rival argument.

This makes disagreement cheaper in a precise sense. Two papers can
disagree while sharing an evidentiary substrate. Each can be traced back
through a source map; neither is entitled to rewrite the cards it finds
inconvenient. The cost is that the archive can no longer hide behind the
apparent coherence of the manuscript. Contradictions remain visible
after the prose has resolved them rhetorically.

\section{5. Typed absence: the ghost can steer but cannot
testify}\label{typed-absence-the-ghost-can-steer-but-cannot-testify}

The most distinctive object in the system may be one that contains no
content at all.

A card may declare \texttt{{[}{[}an\ address{]}{]}} that cannot be
conservatively resolved to another card, platform, concept, or explicit
external object. SLIPCASE preserves the unresolved address as a
\textbf{ghost} instead of deleting it or guessing a fuzzy match.

The ghost changes the meaning of absence. It is not automatically a
broken link. It may be a concept that several branches independently
need but nobody has yet written; an unstable synonym; a granularity
mismatch; a copied assumption; or a genuine missing research object.
Z-SLIP-002 proposes that repeated inbound demand can make such an
absence a conceptual attractor. Z-SLIP-017 adds the necessary epistemic
type rule: a ghost may authorize \emph{search} but may not authorize
\emph{belief}.

That distinction can be written as a small permission table:

\begin{Shaded}
\begin{Highlighting}[]
\NormalTok{GHOST}
\NormalTok{  may prioritize search       yes}
\NormalTok{  may generate a question     yes}
\NormalTok{  may support a claim         no}
\end{Highlighting}
\end{Shaded}

The idea matters beyond ghosts. Research systems often collapse all
information into one scale of confidence. But different objects
authorize different operations. A verified quotation may support a
textual claim. A high-degree unresolved address may justify spending
retrieval effort. A prompt may explain why the system searched a
particular source. A mechanical checksum may establish byte identity.
None is interchangeable with the others.

This typed view of epistemic authority also prevents a subtle form of
circularity. If the archive creates a concept, notices that many of its
own cards refer to the concept, and then treats graph centrality as
evidence that the concept is important in the world, it has converted
self-reference into pseudo-evidence. Ghosts must therefore remain
questions until independent evidence arrives.

\section{6. Two provenance graphs, not
one}\label{two-provenance-graphs-not-one}

The same typing problem appears in AI disclosure.

A prompt can be causally consequential to a scholarly argument without
being evidence for that argument. Z-SLIP-016 states the distinction
compactly:

\begin{Shaded}
\begin{Highlighting}[]
\NormalTok{CAUSES(claim, prompt) != SUPPORTS(claim, prompt)}
\end{Highlighting}
\end{Shaded}

A research artifact therefore has at least two provenance structures.

The \textbf{evidence trace} answers: what supports this substantive
statement?

\begin{Shaded}
\begin{Highlighting}[]
\NormalTok{claim {-}\textgreater{} zettel {-}\textgreater{} source {-}\textgreater{} citekey}
\end{Highlighting}
\end{Shaded}

The \textbf{production trace} answers: what operations and interventions
caused this artifact to take its present form?

\begin{Shaded}
\begin{Highlighting}[]
\NormalTok{artifact {-}\textgreater{} operation {-}\textgreater{} prompt / model / human intervention}
\end{Highlighting}
\end{Shaded}

W3C PROV is intentionally general enough to describe entities,
activities, and agents in a production process {[}2{]}.
Workflow-oriented RO-Crate profiles similarly treat inputs, outputs,
code, and execution provenance as packageable research context {[}4{]}.
The problem here is not the absence of provenance vocabularies. It is
the temptation to let production history stand in for epistemic support.
``The model generated this after reading source X'' is not the same
claim as ``source X warrants this sentence.''

This suggests a better form of AI disclosure. Rather than placing a
binary label on the final artifact and expecting the label to explain
authorship, the system should preserve consequential interventions while
keeping them separate from evidence. The public paper can carry a quiet
identification; the package carries the fuller making history.

The unresolved problem is selection. If every keystroke and prompt is
preserved, provenance becomes noise. If only a few interventions are
kept, the system requires a theory of \textbf{consequential control}:
which operation changed the source set, interpretation, graph
resolution, title, argument, or verification state enough to deserve a
durable receipt?

\section{7. The checkpoint as scholarly
form}\label{the-checkpoint-as-scholarly-form}

The final layer is the checkpoint itself.

BagIt demonstrates a useful preservation principle: a package can have
enough structure to support reliable transfer without understanding the
internal semantics of every payload file {[}1{]}. RO-Crate extends the
research-object idea by aggregating research data and contextual
metadata in a lightweight, machine-readable package {[}3{]}. Git
demonstrates another complementary property: content-addressable storage
can separate an object's identity from its human-facing pathname
{[}5{]}. SLIPCASE borrows from all three tendencies while pursuing a
different object: a research field whose unfinished structure remains
operable.

The root-level \texttt{.txt} cards are therefore not merely a retro
preference for simple files. They encode a recoverability hierarchy.
Z-SLIP-006 summarizes it as: root holds what would be intellectually
grieved; derived folders hold what should be regenerable. A graph,
reader, bibliography view, or paper can be useful and beautifully made
while remaining secondary. If it disappears, the preserved research
objects should be sufficient to produce another one.

The single-file \texttt{index.html} has the opposite role. It is
intentionally monolithic because transmission and canonical storage
optimize for different things. The canonical package favors inspectable
separate files. The sendable capsule favors a recipient who receives one
file and needs immediate access to the cards, graph, paper, prompts, and
reconstruction path. Z-SLIP-012 names this distinction: a sendable
research object is stronger than a shareable interface when it carries
enough state to reconstruct the environment behind the view.

This is the point at which ``research preservation'' becomes ``research
continuation.'' The checkpoint does not merely answer \emph{what
happened?} It should answer:

\begin{itemize}
\tightlist
\item
  What evidence is here?
\item
  What claims are derived from it?
\item
  What remains unresolved?
\item
  Which relations are literal and which are compiled?
\item
  What source or test should be pursued next?
\item
  Which parts of the present interpretation may be discarded safely?
\item
  What can be mechanically verified?
\item
  What could not be verified?
\item
  How does another context rejoin the work?
\end{itemize}

A checkpoint is successful when those questions can be answered without
the originating chat.

\section{8. The model judges; the tool
counts}\label{the-model-judges-the-tool-counts}

A final principle protects the architecture from a mundane but pervasive
failure. Language models are good at interpreting fuzzy material and bad
at pretending that remembered quantities are measurements. SLIPCASE
therefore assigns deterministic assertions to deterministic mechanisms
wherever possible.

Hashes, file counts, relation extraction, citekey inclusion, duplicate
detection, PDF status, and ZIP integrity should be computed.
Interpretation remains a model or human judgment. Z-SLIP-004 describes
this not as ``use tools'' but as a division of epistemic authority.

The distinction needs one qualification. A script can count a badly
defined category perfectly. Determinism proves neither the category nor
the interpretation; it proves only the operation it actually performed.
A strong checkpoint therefore reports several kinds of completeness
separately. Z-SLIP-005 distinguishes source/context coverage,
classification coverage, and mechanical reconciliation. There is no
honest scalar \texttt{COMPLETE:\ YES} when the denominators are
different or unknown.

This is where the checkpoint becomes more than packaging. It publishes
the boundary of what the research process knows about itself.

\section{9. Continuability as a design
criterion}\label{continuability-as-a-design-criterion}

The combined system can be described as a research metabolism, but the
metaphor should not be allowed to imply autonomy or life. Its components
are designed operations: source selection, admission, steering, relation
resolution, compilation, and publication. The value of the metaphor is
narrower. The system is stable not because its interpretations remain
fixed but because it can undergo another transformation without
consuming the evidence of the previous one.

That yields a candidate design criterion for AI-mediated scholarship:

\begin{quote}
A research state is continuable when another competent researcher or
model can recover a warranted next operation, execute it against
identifiable evidence, append the result without rewriting prior
evidence, and reconstruct the relation between the new state and the old
one.
\end{quote}

This criterion is stronger than readability and weaker than full
reproducibility. A recipient may be unable to reproduce every historical
model output and still be able to continue the inquiry responsibly.
Conversely, a recipient may reproduce a PDF byte-for-byte while having
no idea what question remained open.

Continuability therefore connects preservation, provenance, and inquiry
without collapsing them. It asks not whether the archive can replay the
past perfectly, but whether it preserves enough structure for the future
to disagree with the past intelligently.

\section{10. Limits and unresolved
territory}\label{limits-and-unresolved-territory}

Several claims in this architecture remain proposals rather than
established results.

First, no evidence yet shows which zettel fields are necessary or
sufficient for successful continuation. The decisive test is empirical:
give stripped-down research objects to unfamiliar researchers and
measure which operations they can reconstruct.

Second, graph steering may amplify the vocabulary of the archive.
High-degree ghosts and recurring sources are useful signals, but they
can also reward copied assumptions or premature naming convergence.

Third, immutable payloads do not make evidence neutral. Admission and
segmentation are interpretive acts whose histories may themselves
require preservation.

Fourth, the line between reconstructible orchestration state and
irrecoverable research history is unsettled. If the daemon's choice not
to pursue a branch materially shaped the field, the absence of that
decision from the cards may matter.

Fifth, dual provenance needs a practical rule for consequentiality.
Capturing everything is unusable; capturing too little turns making
history into decoration.

Finally, the system has not yet demonstrated long-horizon merges across
genuinely independent checkpoints. Content hashes solve byte identity,
not conceptual identity, synonymy, or intellectual descent.

These are not reasons to close the design. They are the next research
objects the design has managed to expose.

\section{Conclusion}\label{conclusion}

The durable scholarly object is usually imagined as the publication,
dataset, repository, or executable workflow. AI-mediated research
reveals another possibility: the thing worth preserving may be a state
of inquiry that remains capable of becoming different.

The architecture developed here separates atomic evidence-bearing
objects from the controller that allocates attention; separates literal
evidence from compiled field structure; separates support provenance
from production provenance; permits unresolved objects to steer search
without becoming evidence; and treats the paper as a revisable wager
over an append-only history.

The resulting checkpoint is not valuable because it captures everything.
It is valuable because it can state what it has, what it lacks, what it
currently thinks, and what operation could make it think differently.

A readable archive remembers. A continuable archive can be disagreed
with, extended, and made new without losing the path by which it
arrived.

\section{References}\label{references}

{[}1{]} J. Kunze, J. Littman, E. Madden, J. Scancella, and C. Adams.
\textbf{The BagIt File Packaging Format (V1.0)}. RFC 8493, 2018. DOI:
10.17487/RFC8493.

{[}2{]} L. Moreau and P. Missier, eds.~\textbf{PROV-DM: The PROV Data
Model}. W3C Recommendation, 30 April 2013.

{[}3{]} S. Soiland-Reyes et al.~\textbf{Packaging research artefacts
with RO-Crate}. \emph{Data Science} 5(2):97-138, 2022. DOI:
10.3233/DS-210053.

{[}4{]} S. Leo et al.~\textbf{Recording provenance of workflow runs with
RO-Crate}. arXiv:2312.07852, 2023.

{[}5{]} S. Chacon and B. Straub. \textbf{Pro Git}, 2nd ed., section
``Git Internals - Plumbing and Porcelain.'' Git documentation describes
Git as fundamentally a content-addressable filesystem with a
version-control interface layered above it.

\section{Appendix A. Assembly
instrument}\label{appendix-a.-assembly-instrument}

The exact recursive-forage and SLIPCASE assembly prompts used for this
checkpoint are preserved under \texttt{\_PROMPTS/} and embedded in
\texttt{index.html}. Their byte hashes are recorded in
\texttt{\_SLIPCASE/MANIFEST.json}.

\section{Appendix B. Making history}\label{appendix-b.-making-history}

This paper was assembled from the twenty \texttt{Z-SLIP-*} zettels
produced in the working conversation, the supplied prompt architecture,
and a small set of external standards used only to position the design
against existing approaches to provenance and research-object packaging.
The zettel payloads were preserved as cards. The organization and
connective argument of this paper are derived interpretation.

See \texttt{000\_\_MAKING\_HISTORY.txt} for the package-level record.

\section{Appendix C. Replication
path}\label{appendix-c.-replication-path}

Open \texttt{index.html} for the self-contained reader. Open
\texttt{000\_\_RETURN\_PATH.txt} and \texttt{000\_\_REBUILD.txt} for the
checkpoint identity and reconstruction procedure. The canonical evidence
objects are the root zettel cards and preserved source/prompt resources,
not this paper.

\end{document}
